% v2 -- rewritten
% Seções:
%   8.1 - Laboratório virtual de periodontia com SUS-Twin
%   8.2 - Simulação de casos clínicos com DTs
%   8.3 - Avaliação objetiva com DTs (OSCE digital)
%   8.4 - Projeto final: criar um módulo de ensino com DT
\destaque{Nivel-alvo N1 -- Implementacao Guiada: o estudante constroi competencias praticas em periodontia utilizando o Framework SUS-Twin como metodologia pedagogica.}
\label{ch:educacao}

O ensino da periodontia baseia-se tradicionalmente na observação clínica
supervisionada. Com o SUS-Twin, é possível inverter essa lógica: o estudante
manipula Gêmeos Digitais\footnote{Um Gêmeo Digital (DT) é uma réplica
computacional de um paciente real, construída a partir de dados clínicos,
imagens e exames, que permite simular procedimentos sem risco ao paciente.}
de pacientes periodontais antes do contato com o caso real, acelerando o
raciocínio clínico e reduzindo a curva de aprendizado. Este capítulo
apresenta roteiros práticos para quatro encontros laboratoriais, todos
suportados pelo ecossistema SUS-Twin e pela plataforma OpenCode.

% =============================================================================
\section{Laboratório virtual de periodontia com SUS-Twin}
\label{sec:lab-virtual}

O laboratório virtual\footnote{Ambiente computacional que replica a
experiência de um laboratório didático, executado inteiramente em software,
sem necessidade de equipamentos físicos odontológicos.} é o ponto de partida.
Cada estação de trabalho requer Python 3.11, Jupyter Notebook\footnote{
Plataforma interativa que combina código executável, visualizações e texto
explicativo em um único documento, muito usada no ensino de programação
científica.}, PySUS (consulta a dados do DATASUS) e Open3D (visualização 3D
de tomografias). O arquivo abaixo é um modelo inicial de caderno Jupyter que
carrega um estudo de caso em formato DT\footnote{Um estudo de caso em formato
DT é um arquivo JSON contendo dados clínicos, exames de imagem segmentados e
metadados do paciente, prontos para serem manipulados programaticamente.}:

\begin{lstlisting}[language=Python,
caption={modelo\_dt\_starter.ipynb -- carrega e inspeciona um DT
periodontal}]
# Jupyter Notebook: Periodontia com SUS-Twin
import pysus
import open3d as o3d
import json, numpy as np

# Carrega o caso clinico no formato DT
with open("casos/paciente_001.json") as f:
    caso = json.load(f)

print(f"Paciente: {caso['id']} | Idade: {caso['idade']}")
print(f"CID: {caso['cid']} | Sextantes afetados: {
      len(caso['sextantes_afetados'])}")

# Converte a CBCT segmentada para nuvem de pontos 3D
nuvem = o3d.geometry.PointCloud()
nuvem.points = o3d.utility.Vector3dVector(
    np.array(caso['cbct_pontos']))
o3d.visualization.draw_geometries([nuvem],
    window_name="DT Periodontal -- SUS-Twin")
\end{lstlisting}

O código acima carrega um DT, exibe seus metadados clínicos e renderiza a
tomografia computadorizada em 3D. O estudante deve executar célula a célula,
explorando os atributos do caso. O SUS-Twin fornece uma biblioteca com 20
casos periodontais reais anonimizados para este fim, todos acessíveis via
repositório do Framework SUS-Twin.

% =============================================================================
\section{Simulação de casos clínicos com DTs}
\label{sec:simulacao-casos}

Cada caso a seguir foi projetado para uma sessão laboratorial de 2 horas. A
progressão de dificuldade é intencional (A $<$ B $<$ C), sempre no contexto do
Framework SUS-Twin.

\subsection*{Caso A -- Periodontite crônica generalizada (CID K05.3)}

O arquivo \texttt{casos/paciente\_A.json} contém dados de sondagem
periodontal (SIA) de um paciente com periodontite crônica. O estudante deve
identificar quantos sextantes\footnote{Na odontologia, a arcada dentária é
dividida em seis sextantes: três na arcada superior (direito, anterior,
esquerdo) e três na inferior. Cada sextante é avaliado separadamente quanto
à presença de doença periodontal.} estão afetados (perda de inserção
clínica $\geq$ 5~mm em dois ou mais sítios) e classificar a severidade
segundo os critérios da AAP/CDC. Espera-se que o aluno produza um mapa
térmico da perda de inserção utilizando matplotlib e um parágrafo de
diagnóstico. \textbf{Rubrica:} acertar o número de sextantes (2~pts),
classificação correta (2~pts), mapa térmico legível (1~pt).

\subsection*{Caso B -- Recessão gengival + lesão de furca}

O arquivo \texttt{casos/paciente\_B.json} descreve um caso de recessão
gengival classe III de Miller associada a lesão de furca\footnote{Furca (ou
furcação) é a região de bifurcação ou trifurcação das raízes de dentes
multirradiculares (molares). A lesão de furca ocorre quando a doença
periodontal atinge essa região, dificultando o tratamento e exigindo
avaliação tridimensional.} classe II (perda óssea horizontal de 3~mm). O
estudante deve segmentar a CBCT manualmente no Open3D e aplicar o Método dos
Elementos Finitos (FEM) disponível no módulo \texttt{sustwin.fem} para
simular o efeito de um enxerto gengival. A análise esperada inclui a
distribuição de tensões antes e depois da simulação.
\textbf{Rubrica:} segmentação correta (2~pts), FEM executado sem erros
(2~pts), interpretação biomecânica (1~pt).

\subsection*{Caso C -- Avaliação de risco de implante em paciente fumante}

O arquivo \texttt{casos/paciente\_C.json} contém um DT de paciente tabagista
(20 cigarros/dia) com indicação de implante no elemento 36. Sabe-se que o
tabagismo\footnote{O fumo reduz a vascularização gengival, compromete a
resposta inflamatória e diminui a taxa de osseointegração, aumentando em
até 300\% o risco de falha precoce de implantes dentários.} reduz
significativamente a taxa de sucesso de implantes. O estudante deve executar
uma validação cruzada K-Fold (K=5) no modelo preditivo de risco de falha
fornecido pelo módulo \texttt{sustwin.risk\_assessment}, comparando a
acurácia com e sem a variável ``tabagismo''. O relatório final deve incluir
a matriz de confusão e a curva ROC.
\textbf{Rubrica:} K-Fold implementado corretamente (2~pts), matriz de
confusão (1~pt), curva ROC com AUC (1~pt), discussão do impacto do
tabagismo (1~pt).

% =============================================================================
\section{Avaliação objetiva com DTs (OSCE digital)}
\label{sec:osce-digital}

O OSCE\footnote{Objective Structured Clinical Examination: exame clínico
objetivo e estruturado no qual o estudante percorre estações com tarefas
específicas, cada uma avaliada por uma rubrica padronizada.} digital proposto
utiliza quatro estações, todas executadas no ambiente SUS-Twin, com duração
de 25 minutos cada:

\begin{enumerate}
\item \textbf{Segmentação CBCT:} o aluno recebe uma tomagem não processada e
deve segmentar o osso alveolar e as raízes dentárias no Open3D.
\item \textbf{Consulta DATASUS:} utilizando PySUS, o aluno deve extrair a
taxa de mortalidade por doenças periodontais no estado de São Paulo entre
2018--2023, gerando um gráfico de série temporal.
\item \textbf{Simulação FEM:} o aluno executa o modelo de elementos finitos
do Caso B (seção~\ref{sec:simulacao-casos}) com parâmetros alterados
(força oclusal de 200~N em vez de 150~N) e reporta a nova distribuição de
tensões.
\item \textbf{Interpretação de resultados:} o aluno recebe três outputs do
SUS-Twin e deve redigir um parecer clínico de 5 linhas para cada um.
\end{enumerate}

O código abaixo implementa um corretor automático\footnote{A correção
automática (auto-grading) utiliza scripts que comparam a saída gerada pelo
aluno com valores esperados previamente cadastrados pelo professor. O
resultado é imediato e objetivo, embora não substitua a avaliação
qualitativa do raciocínio clínico.} que verifica as respostas objetivas:

\begin{lstlisting}[language=Python,
caption={auto\_grader\_osce.py -- corretor automático do OSCE digital}]
import json, numpy as np

RESPOSTAS_ESPERADAS = {
    "sextantes_afetados": 4,
    "fem_tensao_max_pas": 12.7,
    "auc_risco_implante": 0.89
}

def corrigir(respostas_aluno: dict) -> dict:
    resultado = {}
    for chave, esperado in RESPOSTAS_ESPERADAS.items():
        obtido = respostas_aluno.get(chave)
        if isinstance(esperado, float):
            resultado[chave] = abs(obtido - esperado) < 0.1
        else:
            resultado[chave] = obtido == esperado
    return resultado

notas = corrigir(json.load(open("respostas_aluno.json")))
print(f"Nota objetiva: {sum(notas.values())}/{len(notas)}")

\end{lstlisting}

O corretor pode ser integrado ao ambiente SUS-Twin para devolução imediata
de feedback. O professor mantém a rubrica qualitativa para as questões
discursivas, enquanto a parte técnica é avaliada automaticamente.

% =============================================================================
\section{Projeto final: criar um módulo de ensino com DT}
\label{sec:projeto-final}

Como avaliação somativa da disciplina, cada aluno (ou dupla) deve criar um
módulo de ensino original utilizando o Framework SUS-Twin. O módulo deve
integrar quatro componentes obrigatórios:

\begin{itemize}
\item \textbf{Caso clínico com DT:} selecionar um paciente periodontal,
anonimizá-lo e estruturar os dados no formato DT do SUS-Twin, incluindo
sondagem, CBCT segmentada e histórico médico.
\item \textbf{Consulta PySUS:} associar ao caso uma consulta ao DATASUS que
contextualize epidemiologicamente a condição do paciente (prevalência
regional, faixa etária, comorbidades associadas).
\item \textbf{Visualização 3D:} gerar ao menos duas visualizações com Open3D
(ex.: nuvem de pontos da CBCT e malha FEM) que ilustrem aspectos relevantes
do plano de tratamento.
\item \textbf{Rubrica de validação:} elaborar uma rubrica com no mínimo 5
critérios e 3 níveis de proficiência (Insuficiente, Proficiente, Avançado),
cobrindo desde a corretude técnica até a clareza da exposição.
\end{itemize}

O projeto deve ser entregue como um repositório no GitHub contendo o caderno
Jupyter, os dados do DT, o script de consulta PySUS e a rubrica em PDF. A
avaliação final é conduzida por pares (peer review) utilizando a rubrica
criada pelo próprio grupo, sob supervisão do docente. O ecossistema SUS-Twin
disponibiliza templates e exemplos completos no repositório oficial do
OpenCode para auxiliar a elaboração do módulo, consolidando a metodologia
pedagógica do Framework SUS-Twin.
